\documentclass[11pt]{article}
\usepackage{amsmath,amssymb,amsthm}
\usepackage[margin=1in]{geometry}
\newtheorem{theorem}{Theorem}
\newtheorem{lemma}{Lemma}
\newtheorem{corollary}{Corollary}
\begin{document}
\title{V113: Common-Dominator Dynamic Programming for Minimum-Overlap MUX Return Flows}
\author{NC0\_k-Avoid Laboratory}
\date{}
\maketitle

Fix two signed essential ternary MUX outputs with common selector $v$, and choose first
branches with opposite source phases.  Let their destinations be $d_0,d_1$.  Return
routes end at $v$, and outputs whose selector is $v$ are excluded from the return graph.
Each remaining output is treated as one gate resource.

\paragraph{Common gate dominators.}
A return gate $h$ is common if every $d_0\to v$ route and every $d_1\to v$ route uses
$h$.  Let $H$ be the set of such gates.

\begin{lemma}[Dominator chain]
The gates of $H$ are linearly ordered from the two return sources to the sink.
\end{lemma}

\begin{proof}
Add a super-source joined to $d_0$ and $d_1$.  A gate is in $H$ exactly when its split
gate-node dominates the sink $v$.  Dominators of one vertex are the ancestors of that
vertex in the dominator tree and hence form a chain.
\end{proof}

\begin{theorem}[Minimum overlap]
For the fixed first-branch pair,
\[
 \min\{|P_0\cap P_1|:P_i\text{ is a }d_i\to v\text{ return route}\}=|H|,
\]
where intersection counts shared output-gate resources.
\end{theorem}

\begin{proof}
Every route from either source uses every gate in $H$, giving the lower bound.

For the upper bound, give each gate of $H$ capacity two and every other gate capacity
one; transition arcs and variables have capacity two.  A super-source sends one unit to
each $d_i$.  If a cut of capacity one existed, it could not be a common gate, whose
capacity is two, and one source arc alone cannot separate the other source.  Thus the
unit cut would be a noncommon gate whose deletion separates every super-source-to-$v$
route.  That gate would lie on every route from both sources, contradicting its being
noncommon.  Max-flow/min-cut and integrality therefore give a two-flow whose two paths
share no noncommon gate.  They share exactly $H$.
\end{proof}

\begin{corollary}[Dominator intervals]
Every minimum-overlap pair shares exactly $H$ and is gate-disjoint between consecutive
common dominators.  Conversely, gate-disjoint route pieces in each interval stitch to a
minimum-overlap pair.
\end{corollary}

At a common MUX gate $h_j$, the only shared branch state is
\[
 (b_0,b_1)\in\{0,1\}^2.
\]
Thus every dominator has four states.

\begin{lemma}[Polynomial phase transition]
For a fixed state at $h_j$ and fixed next state at $h_{j+1}$, feasibility of a
minimum-overlap target-compatible transition is decidable in polynomial time.
\end{lemma}

\begin{proof}
The target required at $h_j$ by one route depends only on its branch at $h_j$ and the
source phase of the first selected output branch in the following interval.  Enumerate
the first gate and branch for each route.  Reject pairs that reuse one noncommon gate or
induce different target bits at $h_j$.  Once these first choices are fixed, complete
the two tails to the next boundary with a gate-capacitated integral two-flow having
capacity one on every noncommon gate.  There are only polynomially many first-choice
pairs and an ordinary max-flow test for each.
\end{proof}

\begin{theorem}[Optimum-face completeness]
For a fixed opposite-phase first-branch pair, there is a deterministic polynomial-time
algorithm deciding whether some minimum-overlap return pair is target-compatible on all
shared gates.  When one exists, the algorithm constructs a missing output word.
\end{theorem}

\begin{proof}
Run dynamic programming along the ordered common dominators.  The state at one common
gate is its pair of branch choices, so there are four states.  Use the previous lemma to
create transitions between states in consecutive intervals.

Soundness follows because every accepted transition uses gate-disjoint noncommon
resources, while the two traversals of each common gate require the same target bit.
The resulting two opposite-phase cycles therefore coexist in one fixed-output 2-CNF and
force opposite values of the root selector, so the completed target word is missing.

For completeness, take any target-compatible minimum-overlap pair.  By the minimum
overlap theorem it shares exactly $H$.  Its branch pair at each common gate is one of
the four DP states.  In every interval its first branches occur in the explicit
enumeration and its remaining gate-disjoint tails witness the max-flow feasibility
check.  Hence all of its state transitions are present and the DP accepts.
\end{proof}

\begin{corollary}[Polynomial scan]
Scanning every pair of outputs with a repeated selector and every opposite-phase pair of
first branches remains polynomial.  Therefore V113 finds a missing word whenever any
such first pair has a target-compatible minimum-overlap return pair.
\end{corollary}

\begin{theorem}[Nonserial strict family]
For every $d\ge1$ there is a signed MUX family with
\[
 n=7(d+1),\qquad m=n+1,
\]
whose distinguished return pair has minimum overlap $d$, while the same optimum face
contains both a target-compatible and a target-incompatible pair.  The support uses
three-variable lobes and is outside the exact two-variable-lobe class recognized by
V112.
\end{theorem}

\begin{proof}
Use the periodic $k=3$ signing recorded in the V113 reference implementation.  The
$d$ hub gates are on every return route, so they are common dominators.  Explicit left
and right lobe routes are otherwise gate-disjoint and therefore attain overlap $d$.
For one explicit route pair the two target requirements agree at every shared hub; for
a second explicit pair they disagree at every repeated block.  Both pairs are overlap
optimal, but only the first is target-compatible.
\end{proof}

\begin{theorem}[Exact V102 backdoor]
For the strict family,
\[
 \beta_{\mathrm{V102}}=5(d+1)=\frac{5n}{7}.
\]
\end{theorem}

\begin{proof}
For one three-variable lobe $x_0,x_1,x_2$ with exit hub $z$, the MUX backdoor conditions
are cyclically
\[
 x_i\in B\quad\text{or}\quad (x_{i+1}\in B\text{ and }z\in B).
\]
If $z\in B$, the conditions reduce to a vertex cover of a triangle and require two lobe
variables.  If $z\notin B$, all three lobe variables are forced.  Two lobes therefore
cost at least four selected lobe variables when $z$ is selected and six otherwise;
including $z$ gives layer cost at least five or six respectively.  Summing over $d+1$
layers gives $\beta\ge5(d+1)$.  Selecting every exit hub and any two variables from each
lobe attains equality and also satisfies the central/shared hub gates.
\end{proof}

\paragraph{Boundary.}
V113 solves the zero-excess regime
\[
 \Delta=\operatorname{minCompatibleOverlap}-\operatorname{minOverlap}=0
\]
for every fixed opposite-phase first pair.  It does not prove that every MUX circuit has
a compatible optimum.  The case $\Delta>0$, all essential MUX/bijunctive circuits,
unrestricted $NC^0_3$-Avoid, general circuit lower bounds, and P versus NP remain open.
Novelty and priority are not established.

\end{document}
